% Source: https://artofproblemsolving.com/community/c3h3601683__math_problem_inspired_by_mathcounts_
% Stable identity: aops-topic-3601683-q01
Any parenthesization makes each number a factor with exponent $+1$ or $-1$. The first number has exponent $+1$ and the second exponent $-1$. For each prime, changing signs alters its total exponent by an even number. The total exponents of $2,3,5,7$ are all odd: respectively $1+14$, $-1+6$, $3$, and $1$. If the value is an integer, each of these exponents is a nonnegative odd integer, hence at least one. Thus every integer value is divisible by $2\cdot3\cdot5\cdot7=210$. This bound is attained by\[4\div\bigl((6\div8\div9\div10\div12\div14)\div(15\div16\div18\div20)\bigr)=210,\]where unparenthesized chains inside each group are evaluated left to right. Therefore the answer is $\boxed{210}$.
